\section{CENÁRIO}

A presença constante dos celulares no dia a dia e o acesso à internet em tempo real mudaram a forma como as pessoas se comunicam, tornando a rapidez um fator decisivo para resolver problemas nas cidades. Nesse contexto, as tecnologias digitais facilitam a criação de redes de apoio comunitário, onde os próprios moradores participam compartilhando informações e ajudando em situações de interesse coletivo (CASTELLS, 2020). No caso da proteção aos animais domésticos, essa união de esforços é essencial para organizar ações rápidas de resgate e acolhimento entre pessoas que moram ou circulam na mesma região.

Com os smartphones consolidados como a principal forma de acesso à internet no Brasil (CNDL \& SPC BRASIL, 2024), os pedidos de ajuda comunitária passaram naturalmente para o ambiente virtual. No entanto, apenas publicar pedidos de socorro em redes sociais comuns não resolve a urgência de quem perdeu um animal. Aplicativos criados para essa finalidade precisam ser fáceis de usar e seguros: devem organizar fotos e características do animal como espécie, porte e cor, mostrar as ocorrências em um mapa por proximidade e garantir a privacidade das conversas, respeitando as regras da Lei Geral de Proteção de Dados Pessoais (BRASIL, 2018).

\subsection{O Crescimento da População de Animais Domésticos no Brasil}

A importância dos animais de estimação na vida dos brasileiros é comprovada por pesquisas oficiais. Segundo dados do Instituto Brasileiro de Geografia e Estatística (IBGE, 2020), 46,1\% dos lares brasileiros tinham pelo menos um cão e 19,3\% contavam com pelo menos um gato, somando dezenas de milhões de animais que fazem parte das famílias. Dados do Conselho Federal de Medicina Veterinária indicam que essa população já passa de 70 milhões de animais, mostrando que a posse responsável e os cuidados com o bem-estar animal são desafios constantes para as cidades (CFMV, 2015).

Do ponto de vista da lei, o resgate e o cuidado com animais encontrados encontram apoio no Código Civil (BRASIL, 2002), na Lei Geral de Proteção de Dados (BRASIL, 2018) e nas leis de proteção contra maus-tratos e defesa dos animais (BRASIL, 2020; BRASIL, 2026). Essas leis reforçam que as plataformas de divulgação devem trabalhar com responsabilidade, protegendo os dados de contato dos usuários e garantindo a veracidade das informações cadastradas.

\subsection{A Problemática da Fuga e Desaparecimento de Animais}

Quando um cão ou gato se perde em uma cidade, as primeiras horas são as mais importantes para o resgate. Sem saber para onde ir, o animal caminha sem rumo, correndo risco constante de atropelamento, fome, chuva e agressões físicas. Em momentos de desastres e enchentes, como as que atingiram o estado do Rio Grande do Sul em 2024, a situação fica ainda mais grave, com milhares de animais separados dos donos e uma sobrecarga enorme em abrigos e protetores voluntários (FAROL, 2020; UOL, 2024).

Diariamente, tutores enfrentam a aflição de ter seus animais sumidos, sofrendo com a desorganização das informações espalhadas pela internet e com a falta de canais que avisem quem está por perto. Quando não existe um local centralizado para cadastrar a ocorrência, o esforço da comunidade se perde, consome horas preciosas logo após a fuga e diminui muito a chance de o animal voltar para casa.

\subsection{O Papel dos Smartphones e as Limitações dos Meios Tradicionais e Redes Sociais}

Apesar da importância do tema, os meios tradicionais de busca funcionam de forma limitada. Fazer e colar cartazes em papel nos postes gasta tempo de impressão, custa dinheiro e tem pouco alcance, além de os papéis rasgarem ou se estragarem na primeira chuva. 

Por outro lado, as redes sociais e os grupos de mensagens sofrem com a rapidez com que as postagens somem: as mensagens de desaparecimento se perdem no meio de propagandas e outros posts, não possuem busca por mapa e obrigam os tutores a divulgar números de telefone pessoais na internet, gerando medo de golpes e trotes.

Para criar uma ferramenta que resolva de verdade esses problemas, foi preciso entender diretamente o que as pessoas passam no momento da busca. O Capítulo 3 apresenta a metodologia adotada no trabalho, trazendo a pesquisa prática realizada com a comunidade e o estudo detalhado dos principais sistemas similares disponíveis no mercado.
\clearpage
